\chapter{Introduction}\label{chap:introduction}

Tout matériau réel est imparfait : un atome manque, un autre occupe la mauvaise place ou une impureté s'est immiscée. À l'échelle macroscopique, ces accidents se comptent par milliards et déterminent souvent ce que le matériau fait de mieux comme de pire : ils limitent la mobilité des électrons dans un transistor, déterminent la couleur de certains cristaux, ancrent les vortex d'un supraconducteur ou font apparaître un moment magnétique là ou le matériau parfait n'en possède pas. Comprendre ces défauts, leur nature, leur concentration et surtout leur influence sur les propriétés du matériau, constitue un enjeu central de la science des matériaux.


\section{Contexte et motivation}
\label{sec:contexte}

\textcolor{red}{À reformuler}


Cette question se pose avec une acuité particulière lorsque le
matériau est réduit à une seule couche d'atomes. Dans un matériau
bidimensionnel, chaque atome est un atome de surface : un défaut n'y
est jamais enfoui dans un volume qui en dilue l'effet, il est exposé,
et son influence sur les électrons et les vibrations du cristal l'est
tout autant. C'est ce qui fait de ces matériaux, et du plus connu
d'entre eux, le graphène, le terrain idéal pour étudier les défauts à
l'échelle atomique.

\bloc{Le graphène}

Le graphène est une feuille de carbone d'un seul atome d'épaisseur, dont les atomes sont arrangés sur un réseau en nid d'abeille. Isolé pour la première fois en 2004 par l'exfoliation mécanique du
graphite~\cite{novoselovElectricFieldEffect2004}, il est à la fois l'un des matériaux les plus simples que l'on puisse imaginer et l'un des plus singuliers~\cite{geimRiseGraphene2007}. Sa cellule unitaire ne contient que deux atomes. Pourtant, avec un module de Young de l'ordre du térapascal et une résistance à la rupture de $\sim 130$~GPa, la plus haute jamais mesurée~\cite{leeMeasurementElasticProperties2008}, c'est le matériau le plus résistant que l'on connaisse. Il conduit la chaleur mieux que le cuivre, à $\sim 5 \times 10^{3}$~W/mK~\cite{balandinSuperiorThermalConductivity2008}, et les électrons avec une mobilité qui dépasse $10^{5}$~cm$^{2}$\,V$^{-1}$\,s$^{-1}$~\cite{bolotinUltrahighElectronMobility2008}. En outre, ce matériau, qui bat tous les records, est transparent à $97.7$~\% dans le visible~\cite{nairFineStructureConstant2008}.

Sa richesse ne tient pas à la complexité de sa chimie, mais à sa géométrie et à la nature des liens entre ses atomes. Le réseau en nid d'abeille mène à une structure électronique particulière : la bande de valence $\pi$ touche linéairement la bande de conduction $\pi^*$ aux points de Dirac $K$ et $K'$. À ces points, le gap et la densité d'états s'annulent : le graphène est à l'intersection d'un métal et d'un semi-conducteur, d'où sa catégorisation en temps que \textit{semi-métal}. Ses électrons de basse énergie se comportent comme des particules relativistes sans masse~\cite{castronetoElectronicPropertiesGraphene2009}. Cette structure électronique est à l'origine d'une grande partie de ses propriétés emblématiques et, nous le verrons, de la richesse de sa réponse à la lumière.

L'introduction de défauts vient perturber cette géométrie particulière et peut modifier à la fois les états électroniques, les vibrations du réseau et leur interaction. Le graphène offre ainsi un système où il est possible de relier de manière directe une modification locale de sa structure à ses conséquences observables, pourvu que nous disposons d'une sonde pour mesurer l'écart.
\textcolor{red}{1: propriétés connues donc plus facile de quantifier l'effet des défauts. 2: systèmes singuliers plus difficle à traiter, meilleur test pour les méthodes.}\textcolor{red}{Inverser bloc Raman et défauts}

\bloc{La spectroscopie Raman}

La lumière est une sonde de choix : éclairer un cristal avec un laser et regarder la lumière qui en ressort est, en apparence, la plus simple des expériences. Presque toute la lumière diffusée a la couleur du laser sauf une infime fraction qui décalée en fréquence : elle a laissé dans le matériau un quantum de vibration, un phonon, et ce décalage en mesure l'énergie. C'est l'effet Raman~\cite{ramanNewTypeSecondary1928}, et le spectre de ces décalages est une véritable empreinte digitale du réseau cristallin du matériau.

Pour le graphène, cette empreinte est devenue un outil de caratérisation incontournable~\cite{ferrariRamanSpectrumGraphene2006a,ferrariRamanSpectroscopyVersatile2013a}. La mesure est rapide, non destructive, se fait à l'air ambient sans aucune préparation complexe de l'échantillon et permet, sur un seul spectre, d'en apprendre sur le nombre de couches, la déformation (\textit{strain})~\cite{mohiuddinUniaxialStrainGraphene2009}, le dopage et le désordre. Le spectre du graphène parfait est dominé par deux pics : le pic G et le pic 2D. Le premier est la signature ordinaire d'un phonon dans le matériau, tandis que le second ne l'est pas : il fait intervenir deux phonons à la fois et, dans la plupart des matériaux, un tel processus est trop faible pour être observé. Dans le graphène, il est pourtant souvent plus intense que le pic G. La raison tient à la structure électronique décrite précédemment : en l'absence de gap, tout laser visible peut induire des transitions électronique réelles, ce qui amplifie fortemement la diffusion. Nous parlons alors de diffusion Raman \textit{résonante}.

Cette résonance est à la fois ce qui contribue à la richesse du spectre Raman du graphène et ce qui complique son calcul théorique. Pour décrire correctement le spectre résonant, il faut considérer le processus à l'échelle de ses constituants : les quanta de lumière, les électrons et les phonons, ainsi que les interactions entre eux. À cette échelle, une description quantique explicite du processus est nécessaire. \textcolor{red}{faire un pont avec les défauts : pour quantifier l'effets des défauts sur les des particules microscopiques il faut une théorie quantique.}

\bloc{Les défauts}

Les défauts ne sont pas qu'une curiosité théorique : ils sont dans tous les échantillons, qu'on le veuille ou non. La croissance chimique en phase vapeur, le transfert d'un substrat à l'autre, l'exposition à un faisceau d'ions ou simplement l'air ambient en laissent dans toute feuille de graphène~\cite{banhartStructuralDefectsGraphene2011}. Les défauts sont aussi très souvent voulus, car ils peuvent conférer des propriétés additionnelles : greffer de l'hydrogène ou du fluor au graphène lui ouvre un gap~\cite{eliasControlGraphenesProperties2009}, des lacunes ou des atomes étrangers offrent des sites d'ancrage à des molécules ou des atomes métalliques ce qui en fait un capteur~\cite{schedinDetectionIndividualGas2007} et des pores de quelques atomes de diamètre en font une membrane capable de séquencer l'ADN qui la traverse~\cite{garajGrapheneSubnanometreTranselectrode2010}.

Ces défauts ont des signatures dans le spectre Raman : Tuinstra et Koenig ont observé sur le spectre Raman du graphite un pic dont l'intensité croît avec le désordre~\cite{tuinstraRamanSpectrumGraphite1970}, le pic D, absent du graphite parfait. Son intensité relative au pic G, $I_D/I_G$, est depuis devenue la mesure de routine de la densité de défauts dans le graphène~\cite{luccheseQuantifyingIoninducedDefects2010,martinsferreiraEvolutionRamanSpectra2010}. Cette mesure est utile, certes, mais elle ne signale que la présence de défauts, sans dire \textit{lesquels}. Or, un spectre n'a de sens que s'il est comparé à une référence, et l'expérience ne peut pas la fournir. Ce rôle revient à la théorie : prédire, pour chaque type de défaut et chaque concentration, le spectre Raman attendu, à partir de rien d'autre que la structure du matériau et la nature du défaut. C'est
précisément ce que promettent les calculs à partir des premiers
principes, dits \emph{ab initio} : calculer les propriétés du matériau à partir des principes de la mécanique quantique, sans ajuster les paramètres d'un modèle sur l'expérience. Ce mémoire pose, en partie, les fondations d'un tel calcul.

\bloc{Les calculs à partir des premiers principes}

\textcolor{red}{Retravailler ce bloc}
La théorie du spectre Raman résonant du graphène avec des défauts existe déjà~\cite{venezuelaTheoryDoubleresonantRaman2011}, mais les descriptions existantes reposent généralement sur des modèles paramétrés. De telles approches permettent de reproduire les caractéristiques essentielles du spectre, mais sont moins adaptées lorsqu'il s'agit de prédire comment ce spectre change d'un défaut à l'autre, puisque les paramètres du modèle doivent eux-mêmes être déterminés pour représenter le défaut considérer.

Nous voulons prédire plutôt que de reproduire : nous devons décrire le matériau à partir de sa structure atomique et ses électrons. C'est ce que permet la théorie de la fonctionnelle de la densité (DFT)~\cite{hohenbergInhomogeneousElectronGas1964,kohnSelfConsistentEquationsIncluding1965}, un des outils centraux à ce mémoire. Dans cette approche, autant les propriétés électroniques, que celles vibrationnelles, ainsi que celles qui découlent de l'interaction entre les deux, sont calculées directement de la structure atomique. Placer un défaut dans le matériau en modifiant sa structure et en le traitant avec la DFT permet donc de calculer directement la modification de ces propritétés.

Un spectre Raman résonant se calcule à partir d'un ensemble
d'ingrédients. Dans un modèle, chacun est ajusté séparément et rien ne garantit qu'ils soient compatibles entre eux. Dans le cadre de la DFT, ils sont issus d'un même calcul et donc sont cohérents par construction. Les deux mécanismes qui limitent la durée de vie d'un électron excité, sa diffusion par les phonons et sa diffusion par les défauts, peuvent alors être comparés sur un pied d'égalité, ce que nous ferons au terme de ce mémoire. Pouvons-nous obtenir ces ingrédients à partir des premiers principes pour le graphène avec des défauts, de sorte que le spectre Raman puisse ensuite être calculé sans aucun paramètre ajusté ?

\section{État de l'art et positionnement}
\label{sec:intro-etat-art}

\textcolor{red}{Une ou deux phrases d'introduction}


\bloc{Le spectre Raman résonant à partir des premiers principes}

Pour un cristal parfait, loin de toute résonance, le spectre Raman se calule de façon routinière~\cite{lazzeriFirstPrinciplesCalculationVibrational2003} et la plupart des codes DFT l'implémentent. L'approche employée par ces codes suppose que le laser incident dans l'expérience Raman ne crée aucune excitation réelle dans le matériau, une hypothèse qui est systèmatiquement violée à toutes les énergies de laser dans le graphène en raison de son gap nul.

À l'inverse, les calculs de spectres résonants sont nettement plus rares. Ceux-ci requièrent une théorie quantique et des exigences numériques draconiennes, car l'intensité est dominée par les quelques états électroniques qui satisfont exactement la condition de résonance. Gillet, Giantomassi et Gonze~\cite{gilletFirstprinciplesStudyExcitonic2013a} ont établit le cadre au premier ordre en incluant les effets excitoniques, Reichardt et Wirtz l'ont étendu au graphène~\cite{reichardtInitioCalculation$G$2017,reichardtTheoryResonantRaman2019} avec les électrons, les phonons et le couplage électron-phonons tous issus des premiers principes. Le spectre Raman résonant du graphène \textit{parfait} est donc accessible aujourd'hui \textit{ab initio}.

Néanmoins, celui du graphène avec des défauts ne l'est toujours pas. La théorie complète des processus phonon-défaut à été formulée par Venezuela, Lazzeri et Mauri~\cite{venezuelaTheoryDoubleresonantRaman2011}, qui ont reproduit les pics D et 2D ainsi que leur dipersion avec l'énergie du laser, mais traitent le défaut comme une perturbation locale d'un modèle de liaisons fortes et son couplage aux électrons est une constante, c'est-à-dire un paramètre ajusté. C'est précisement ce genre de paramètre que nous voulons retirer. À notre connaissance, aucun spectre Raman résonant avec des défauts n'a été calculé, pour quelque matériau que ce soit, avec tous ses ingrédients issus des premiers principes. Le maillon manquant est le couplage \textit{ab initio} entre les électrons et le défaut et c'est lui que ce mémoire construit.

\bloc{Le couplage électron-défaut à partir des premiers principes}

Le calcul \textit{ab initio} du couplage entre les électrons et un défaut n'est pas une opération de routine en sciences des matériaux. Aucun code communautaire ne le fait et la littérature regroupe quelques articles basés sur des implémentations spécialisées. Le contraste avec le couplage électron-phonon est frappant : ce dernier dispose d'un cadre théorique établi et validé~\cite{giustinoElectronphononInteractionsFirst2017,giustinoElectronphononInteractionUsing2007a} et les codes qui l'implémentent sont distribués à la communauté depuis près de deux décennies~\cite{ponceEPWElectronphononCoupling2016,leeElectronPhononPhysics2023a,zhouPerturboSoftwarePackage2021,gonzeAbinitProjectImpact2020,bruninPhononlimitedElectronMobility2020}. La raison principale est que la présence d'un défaut brise explicitement la symétrie de translation du système, et que le seul moyen de contourner ce problème est de recourir à des super-cellules couteuses.

Deux approches récentes ont posés les bases et chacune apporte la moitié de la solution. En 2019, Lu, Zhou et Bernardi~\cite{luEfficientInitioCalculations2019,luInitioElectrondefectInteractions2020} ont montré comment calculer le couplage dans une base d'ondes planes à partir du potentiel d'une super-cellule contenant le défaut, puis comment l'interpoler par fonctions de Wannier sur les grilles denses qu'exigent les propriétés de transport ou spectrales, telle l'intensité Raman. Leur formalisme traite tout de fois la diffusion au premier ordre dans le potentiel du défaut, l'approximation de Born, ce qui convient aux impuretés chargées des semi-conducteurs, mais échoue pour la plupart des défauts, lesquels induisent généralement des potentiels très intenses. En 2020, Kaasbjerg~\cite{kaasbjergAtomisticMatrixTheory2020} a au contraire traité la diffusion à tous les ordres et a montré que l'approximation de Born peut se tromper d'un ordre de grandeur pour les défauts de matériaux bidimensionnels. Son calcul est toutefois fait dans une base d'orbitales atomiques localisées, sans interpolation vers les grilles denses. 

Le spectre Raman résonant a besoin des deux à la fois : un traitement exacte du potentiel de défaut \textit{et} une évaluation sur des grilles de plusieurs centaines de milliers de points pour statisfaire ses exigences numériques intransigeantes. Combiner les deux n'est pas immédiat et soulève plusieurs problèmes. D'une part, l'interpolation par fonctions de Wannier exige que la grille de la super-cellule et celle des fonctions de Wannier soient commensurable, ce qui forcerait des super-cellules de plusieurs milliers d'atomes et est complètement intraitable. D'autre part, traiter le couplage à tous les ordres sur les grilles requises par le Raman requiert d'effectuer des opérations sur des objets dont la dimension est de plusieurs centaines de milliers, ce qu'aucune machine ne peut faire. Ce mémoire résout ces deux problèmes et livre, pour une lacune dans le graphène, le couplage électron-défaut \textit{ab initio} interpolé par fonctions de Wannier sur des grilles arbitrairement denses et la diffusion qu'il induit, traitée à tous les ordres dans le potentiel de défaut.

\section{Contributions et plan du mémoire}

Les deux premiers chapitres posent le cadre du mémoire. Le chapitre~\ref{chap:théorie} détaille tous les outils : la théorie de la fonctionnelle de la densité pour les électrons, sa version perturbative pour les phonons, les fonctions de Wannier pour l'interpolation vers les grilles denses qu'exige la résonance et l'Hamiltonien complet du système sur lequel tout le reste est construit. Le chapitre~\ref{chap:Raman} rassemble ces outils pour établir la théorie de la diffusion Raman, de sa description semi-classique aux amplitudes quantiques de premier et deuxième ordre, y compris les processus faisant intervenir des défauts. Il se clôt sur la liste des ingrédients que ces amplitudes requièrent : les énergies des électrons $\varepsilon_{n\kv}$ et des phonons $\omega_{\nu\qv}$, le couplage électron-photon $g^{(\lambda)}$, électron-phonon $g^{(\nu)}$ et électron-défaut $M$, ainsi que les taux d'amortissement $\Gamma^\text{ep}_{n\kv}$ et $\Gamma^\text{ed}_{n\kv}$ et la largeur des phonons $\gamma_{\nu\qv}$. Les deux chapitres suivants calculent tous ces ingrédients, hormis le couplage électron-photon qui relève de la même interpolation de Wannier que le reste et n'est pas évalué ici.

Le chapitre~\ref{chap:défauts} est le c\oe ur du mémoire et contient sa contribution originale : le couplage électron-défaut $M$ et le taux $\Gamma_{n\kv}^\text{ed}$ sont construits à partir des premiers principes. D'abord, nous calculons les éléments de matrice $M$ dans une base d'ondes planes en utilisant un raffinement du potentiel de défaut par remplissage par zéros, exact pour un potentiel localisé, qui découple la taille de la super-cellule de la grille des fonctions de Wannier, ouvrant la voie à l'interpolation par ces dernières. Ensuite, nous employons un schéma local qui permet de traiter la diffusion électron-défaut à tous les ordres sans manipuler des objets de très grandes dimensions. Nous appliquons notre méthode à une lacune dans le graphène, non relaxée, sans polarisation de spin et pour une faible concentration de défaut, l'objectif étant de valider la méthode plutôt que de décrire la lacune réelle. Elle livre trois résultats : un état quasi-lié sous la singularité de van Hove de la bande $\pi$, l'échec de l'approximation de Born qui surestime le taux d'une facteur trois en médiane et diverge aux singularités de van Hove, et un couplage dominé par sa partie non locale qui justifie \textit{a posteriori} le remplissage par zéros.

Le chapitre~\ref{chap:eph} calcule le couplage électron-phonon $g^{(\nu)}$, le taux $\Gamma^\text{ep}_{n\kv}$ et la largeur $\gamma_{\nu\qv}$, avec le code \texttt{EPW}, appliqué au graphène parfait. Il est volontairement plus court : la méthode est déjà établie et le chapitre se concentre sur la validation des résultats. Il se termine par la comparaison des deux mécanismes d'amortissement électronique : phonons et défauts. Nous trouvons qu'ils sont du même ordre de grandeur à température ambiente et pour une concentration de défaut de l'ordre du pourcent, de sorte qu'aucun des deux mécanismes n'est négligeable.

Pour finir, le titre du mémoire commence par \og vers \fg{} et ce n'est pas une faute de frappe : le spectre Raman résonant du graphène avec des défauts, dans toute sa gloire, n'est pas calculé ici. Seuls ses ingrédients le sont et ils sont validés puis livrés sur les grilles denses que le spectre résonant exige. Trois étapes séparent encore ces ingrédients d'un calcul d'intensité en bonne et due forme. D'abord, la DFT ne donne pas des énergies électroniques qui sont des énergies d'excitation et la fonctionnelle que nous utilisons sous-estime le coupalge électron-phonon dans le graphène~\cite{lazzeriImpactElectronelectronCorrelation2008} : les deux requièrent une correction par l'approximation $GW$, qui traite la corrélation électronique au-delà de la DFT en combinant la fonction de Green des électrons $G$ et l'interaction de Coulomb écrantée $W$~\cite{hedinNewMethodCalculating1965}. Ensuite, la super-cellule doit être relaxée pour capturer la reconstruction de la lacune dans le graphène~\cite{el-barbaryStructureEnergeticsVacancy2003}. Enfin, le calcul de super-cellule doit être polarisé en spin, puisque la lacune induit un moment magnétique~\cite{yazyevDefectinducedMagnetismGraphene2007}. Ces trois étapes relèvent du calcul et non de la méthode. Nous les laissons ainsi à études doctorales. La méthode de calcul du couplage électron-défaut \textit{ab initio} et son insertion dans la théorie du Raman résonant est ce que mémoire établit.